\documentclass{article}
\usepackage{fontspec}
\setmainfont{Libertinus Serif}
\usepackage[a4paper, margin=1in]{geometry}
\usepackage{amsmath, amsthm, amssymb, mathtools, enumitem, hyperref}
\usepackage[portuguese]{babel}

\begin{document}

\noindent
Curso de Verão \hfill Prof. Bárbara Seelig Pogorelsky \\
\noindent
Álgebra Linear \hfill \textbf{\Large Soluções Lista 3} \hfill \today
\vspace{0.5mm}
\hrule

\section*{Espaço Dual}

\begin{enumerate}
    \item [\textbf{E1}] Sejam \(f: \mathbb{K}^n \to \mathbb{K}\) funcional linear e \(v_1, \ldots, v_n\) base de \(\mathbb{K}^n\). Escreve \(x = \sum x_i v_i \in \mathbb{K}^n\) e aplica \(f\), então tem-se \(f(x) = f(\sum x_i v_i) = \sum x_i f(v_i)\). Define \(f(v_i) \coloneq \alpha_i\) e obtém-se o resultado.
        Reciprocamente, dados \(x = \sum x_i v_i, y = \sum y_i v_i \in \mathbb{K}^n\), toma \(f(x+y) = f\left(\sum x_i v_i + \sum y_i v_i\right) = f\left(\sum (x_i+y_i) v_i\right) = \sum (x_i+y_i) f(v_i) = \sum x_i \alpha_i + \sum y_i \alpha_i = f(x) + f(y)\). Ademais, \(f(\lambda x) = f\left(\lambda\sum x_i v_i\right) = f\left(\sum(\lambda x_i) v_i\right) = \sum \lambda x_i f(v_i) = \lambda \sum x_i \alpha_i = \lambda f(x)\) para todo \(\lambda \in \mathbb{K}\).\qed
    \item [\textbf{E2}] Sejam \(h(x), p(x), q(x) \in \mathcal{P}_2(\mathbb{R})\) tal que \(h(x) = p(x) + q(x)\). Então \(f_\alpha = (p(x)+q(x)) = f_\alpha(h(x)) = h(\alpha) = p(\alpha) + q(\alpha) = f_\alpha(p(x)) + f_\alpha(q(x))\). Toma \(h(x) = \lambda p(x), \lambda \in \mathbb{R}\) e faz \(f_\alpha(\lambda p(x)) = f_\alpha(h(x)) = h(\alpha) = \lambda p(\alpha) = \lambda f_\alpha(p(x))\).\qed
    \item [\textbf{E3}] Escreve \begin{align}
        c_1 f_1 + c_2 f_2 + c_3 f_3 &= 0, \quad c_1, c_2, c_3 \in \mathbb{R} \\
        (c_1 f_1 + c_2 f_2 + c_3 f_3)(p(x)) &= 0, \quad \forall p(x) \in \mathcal{P}_2(\mathbb{R}) \\
        c_1 f_1(p(x)) + c_2 f_2(p(x)) + c_3 f_3(p(x)) &= 0 \\
        c_1 p(0) + c_2 p(1) + c_3 p(2) &= 0
    \end{align}
    toma \(p(x) = (x-2)(x-3)\) e vem \(c_1 = 0\). Faz \(p(x) = (x-1)(x-3)\) e obtém \(c_2 = 0\). Finalmente, com \(p(x) = (x-1)(x-2)\) tem-se \(c_3 = 0\). Logo, \(f_1, f_2\) e \(f_3\) são linearmente independentes.\qed
    \item [\textbf{E4}]
        \begin{enumerate}
            \item [(a)] Nota que \(\dim \mathcal{P}_2(\mathbb{R}) = 3\), então basta mostrar que \(f_1, f_2, f_3\) é LI. Escreve
            \begin{align}
                c_1 f_1 + c_2 f_2 + c_3 f_3 &= 0, \quad c_1, c_2, c_3 \in \mathbb{R} \\
                (c_1 f_1 + c_2 f_2 + c_3 f_3)(p(x)) &= 0, \quad \forall p(x) \in \mathcal{P}_2(\mathbb{R}) \\
                c_1 f_1(p(x)) + c_2 f_2(p(x)) + c_3 f_3(p(x)) &= 0 \\
                c_1 \int_{0}^{1} p(x) \, dx + c_2 \int_{0}^{2} p(x) \, dx + c_3 \int_{0}^{3} p(x) \, dx &= 0.
            \end{align}
            Toma \(p(x) = 1\), \(p(x) = 2x\) e \(p(x) = 3x^2\) e obtém o sistema linear
            \[\begin{cases}
                c_1 + 2 c_2 + 3 c_3 = 0 \\
                c_1 + 4 c_2 + 9 c_3 = 0 \\
                c_1 + 8 c_2 + 27 c_3 = 0
            \end{cases}\]
            cuja única solução é \(c_1 = c_2 = c_3 = 0\). Logo, \(f_1, f_2, f_3\) são linearmente independentes e formam uma base de \((\mathcal{P}_2(\mathbb{R}))^*\).
            \item [(b)] Seja \(p_1, p_2, p_3\) base de \(\mathcal{P}_2(\mathbb{R})\) da qual \(f_1, f_2, f_3\) é base dual. Então, \(f_i(p_j) = \delta_{ij}\), onde \(\delta_{ij}\) é o delta de Kronecker. Escreve \(p_j(x) = a_j x^2 + b_j x + c_j\) e resolve o sistema
            \[\begin{cases}
                f_1(p_j) = \int_{0}^{1} (a_j x^2 + b_j x + c_j) \, dx = \frac{a_j}{3} + \frac{b_j}{2} + c_j = \delta_{1j} \\
                f_2(p_j) = \int_{0}^{2} (a_j x^2 + b_j x + c_j) \, dx = \frac{8 a_j}{3} + 2 b_j + 2 c_j = \delta_{2j} \\
                f_3(p_j) = \int_{0}^{3} (a_j x^2 + b_j x + c_j) \, dx = 9 a_j + \frac{9 b_j}{2} + 3 c_j = \delta_{3j}
            \end{cases}\]
            para \(j = 1, 2, 3\). Obtém-se \(p_1(x) = 3-5x+\frac{3}{2}x^2\), \(p_2(x) = -\frac{3}{2} + 4x - \frac{3}{2}x^2\) e \(p_3(x) = \frac{1}{3}+x - \frac{1}{2}x^2\).
        \end{enumerate}
    \item [\textbf{E5}] 
        \begin{enumerate}
            \item [(a)] Para \(u=(a,b,c) \in \mathbb{R}^3\), escreve \begin{align}
                f(a,b,c) &= f(\alpha_1 u_1 + \alpha_2 u_2 + \alpha_3 u_3), \quad \alpha_1, \alpha_2, \alpha_3 \in \mathbb{R} \\
                &= \alpha_1 f(u_1) + \alpha_2 f(u_2) + \alpha_3 f(u_3) \\
                &= \alpha_1 - \alpha_2 + 3\alpha_3.
            \end{align}
            \item [(b)] Encontra os coeficientes \(\beta_1, \beta_2, \beta_3\) tais que \((2,3,-1) = \beta_1 u_1 + \beta_2 u_2 + \beta_3 u_3 = \beta_1 (1,0,1) + \beta_2 (0,1,-2) + \beta_3 (-1,1,0)\). Resolve o sistema
            \[\begin{cases}
                \beta_1 - \beta_3 = 2 \\
                \beta_2 + \beta_3 = 3 \\
                \beta_1 - 2\beta_2 = -1
            \end{cases}\]
            e obtém \(\beta_1 = 3\), \(\beta_2 = 2\) e \(\beta_3 = 1\). Logo, \(f(2,3,-1) = \beta_1 f(u_1) + \beta_2 f(u_2) + \beta_3 f(u_3) = 0 + 0 + \beta_3 f(u_3) \neq 0\).
        \end{enumerate}
    \item [\textbf{E6}] Seja \(\beta^* = \{f_1, f_2, f_3\}\) base dual de \(\beta = \{v_1, v_2, v_3\}\), onde \(v_1 = (1,0,-1)\), \(v_2 = (1,1,1)\) e \(v_3 = (2,2,0)\). Para qualquer \((a,b,c) \in \mathbb{C}^3\), tem-se \((a,b,c) = \alpha_1 (1,0,-1) + \alpha_2 (1,1,1) + \alpha_3 (2,2,0)\). Resolve o sistema
        \[\begin{cases}
            \alpha_1 + \alpha_2 + 2\alpha_3 = a \\
            \alpha_2 + 2\alpha_3 = b \\
            -\alpha_1 + \alpha_2 = c
        \end{cases}\]
        e obtém \(f_1(a,b,c) = \alpha_1 = \frac{a+b-2c}{3}\), \(f_2(a,b,c) = \alpha_2 = \frac{a+b+c}{3}\) e \(f_3(a,b,c) = \alpha_3 = \frac{a-2b+c}{6}\).
    \item [\textbf{E7}] Seja \(p_1(x) = a_0 + a_1 x + a_2 x^2\), \(p_2(x) = b_0 + b_1 x + b_2 x^2\) e \(p_3(x) = c_0 + c_1 x + c_2 x^2\) base de \(\mathcal{P}_2(\mathbb{R})\) tal que \(\phi_i(p_j) = \delta_{ij}\), onde \(\phi_1, \phi_2, \phi_3\) é a base dual dada. Resolve o sistema
        \[\begin{cases}
            \phi_1(p_j) = \int_{0}^{1} (a_0 + a_1 + a_2)\, dx = \delta_{1j} \\
            \phi_2(p_j) = (a_0 + a_1 x + a_2 x^2)'(1) = \delta_{2j} \\
            \phi_3(p_j) = a_0 = \delta_{3j}
        \end{cases}\]
        para \(j = 1, 2, 3\). Obtém-se \(p_1(x) = 3x - \frac{3}{2} x^2\), \(p_2(x) = -\frac{1}{2}x + \frac{3}{4}x^2\) e \(p_3(x) = 1 - 2x\).
\end{enumerate}

\end{document}